% answers_short.tex —— 2003 年·辽宁卷 · 简短答案速查
% 由 tools/make_short_answers.py 自 answers.tex 的 \gksolhead/\gkanswer 提取，
% 并与 2003/2003_数学真题_辽宁_答案要点.md 逐题比对后生成，勿手改（改源 MD 后重跑）。

\begin{gknotice}
\textbf{使用说明}：本节只给最终答案便于快速核对，完整推导见解析版。答案为编写者独立推导，\textbf{非官方标准答案}。
\end{gknotice}

\gkpart{选择题}
\begin{center}\begin{tabular}{ccccccccccccc}
\toprule
题号 & 1 & 2 & 3 & 4 & 5 & 6 & 7 & 8 & 9 & 10 & 11 & 12 \\
\midrule
答案 & A & D & B & A & D & C & B & C & B & D & C & A \\
\bottomrule
\end{tabular}\end{center}

\gkpart{填空题}
\begin{itemize}[leftmargin=2.6em]
  \item[\textbf{13.}] $-\dfrac{21}{2}$
  \item[\textbf{14.}] 6、30、10
  \item[\textbf{15.}] 120
  \item[\textbf{16.}] ①④
\end{itemize}

\gkpart{解答题}
\begin{itemize}[leftmargin=2.6em]
  \item[\textbf{17.}] (1) 证明见解析；(2) $\dfrac{2\sqrt{3}}{3}$
  \item[\textbf{18.}] $\varphi = \dfrac{\pi}{2}$，$\omega = \dfrac{2}{3}$ 或 $\omega = 2$
  \item[\textbf{19.}] 当 $a \ge 1$ 时，函数的单调递增区间为 $(0, +\infty)$，无单调递减区间；当 $0 < a < 1$ 时，函数的单调递增区间为 $(0, (1-\sqrt{1-a})^2]$ 和 $[(1+\sqrt{1-a})^2, +\infty)$，单调递减区间为 $[(1-\sqrt{1-a})^2, (1+\sqrt{1-a})^2]$。
  \item[\textbf{20.}] (1) $\xi$ 的概率分布为 $P(\xi=0)=\dfrac{3}{25}$，$P(\xi=1)=\dfrac{2}{5}$，$P(\xi=2)=\dfrac{28}{75}$，$P(\xi=3)=\dfrac{8}{75}$；$\eta$ 的概率分布为 $P(\eta=0)=\dfrac{8}{75}$，$P(\eta=1)=\dfrac{28}{75}$，$P(\eta=2)=\dfrac{2}{5}$，$P(\eta=3)=\dfrac{3}{25}$；(2) $E\xi = \dfrac{22}{15}$，$E\eta = \dfrac{23}{15}$
  \item[\textbf{21.}] (1) 证明见解析；(2) $\left(0, \dfrac{1}{3}\right)$
  \item[\textbf{22.}] 存在。当 $a > \dfrac{\sqrt{2}}{2}$ 时，$E\left(0, \dfrac{a}{2} - \dfrac{\sqrt{4a^2 - 2}}{4}\right)$，$F\left(0, \dfrac{a}{2} + \dfrac{\sqrt{4a^2 - 2}}{4}\right)$，定值为 $a$；当 $a = \dfrac{\sqrt{2}}{2}$ 时，$E, F$ 重合于点 $\left(0, \dfrac{\sqrt{2}}{4}\right)$，定值为 $\dfrac{\sqrt{2}}{2}$；当 $0 < a < \dfrac{\sqrt{2}}{2}$ 时，$E\left(-\dfrac{\sqrt{2 - 4a^2}}{4}, \dfrac{a}{2}\right)$，$F\left(\dfrac{\sqrt{2 - 4a^2}}{4}, \dfrac{a}{2}\right)$，定值为 $\dfrac{\sqrt{2}}{2}$。
\end{itemize}
